\section{The Spectral Attractor}
\label{sec:spectral-attractor}

\subsection{Universal Convergence}

When the Steersman operates with spectral regularization alone---no contrastive
loss specifying which channels should couple preferentially---the bridge
converges to a characteristic connectivity level independent of the number of
channels. We term this the \emph{spectral attractor}.

\begin{table}[h]
\centering
\begin{tabular}{lccc}
\toprule
\textbf{Experiment} & $n$ & \textbf{Final Fiedler} & \textbf{Topology} \\
\midrule
H-ch3  & 3  & 0.0951 & spectral-only \\
H-ch4  & 4  & 0.0918 & spectral-only \\
E-001  & 6  & 0.0836 & emanation (spectral effective) \\
H-ch8  & 8  & 0.0944 & spectral-only \\
H-ch12 & 12 & 0.1019 & spectral-only \\
\midrule
\multicolumn{2}{l}{\textbf{Mean $\pm$ SD}} & $0.0934 \pm 0.0065$ & \\
\multicolumn{2}{l}{\textbf{Band}} & 19.4\% (min/max) & \\
\bottomrule
\end{tabular}
\caption{Spectral attractor convergence across channel counts $n \in \{3, 4, 6, 8, 12\}$.
All experiments converge to Fiedler $\approx 0.09$ regardless of $n$.
E-001 (emanation architecture) converges to the same attractor, confirming
that hierarchical constraints do not override the spectral equilibrium.}
\label{tab:spectral-attractor}
\end{table}

The spectral attractor at Fiedler $\approx 0.09$ is stable across a $4\times$
range of channel counts ($n{=}3$ to $n{=}12$), with only 19.4\% variation
(Figure~\ref{fig:spectral-attractor}).
This invariance suggests the attractor is a property of the spectral
regularization objective itself, not of the bridge dimensionality.

\begin{figure}[t]
  \centering
  \includegraphics[width=0.7\textwidth]{paper4/figures/fig_spectral_attractor.pdf}
  \caption{Spectral attractor: universal Fiedler convergence across channel
  counts $n \in \{3, 4, 8, 12\}$. The shaded band shows 19.4\% variation
  around the mean (0.0934). The spectral attractor is the bridge's
  unprogrammed connectivity equilibrium.}
  \label{fig:spectral-attractor}
\end{figure}

\subsection{Bifurcation Under Contrastive Loss}

Adding the contrastive loss breaks the spectral attractor. At $n{=}6$ with
RD contrastive pairs, the Fiedler drops from $\sim$0.09 (attractor) to
$9.0 \times 10^{-5}$---a factor of 1{,}000$\times$. The contrastive loss
does not merely redirect connectivity; it \emph{concentrates} it along the
specified axes while suppressing it everywhere else.

This bifurcation reveals the two independent control channels of the
Steersman:
\begin{enumerate}
  \item \textbf{Spectral regularization} provides \emph{connectivity}---how
    much total coupling the bridge develops.
  \item \textbf{Contrastive loss} provides \emph{topology}---which channels
    receive the coupling budget.
\end{enumerate}

Without contrastive specification, connectivity distributes uniformly across
all channel pairs, producing the attractor. With specification, the same
connectivity budget concentrates along the specified axes, producing
block-diagonal structure.

\subsection{Emanation as Spectral Attractor Variant}

E-001 uses the emanation architecture (a single master bridge with per-layer
offsets and coherence monitoring) instead of independent per-layer bridges.
Its final state---Fiedler 0.084, co/cross 1.12---places it squarely in the
spectral attractor regime. The hierarchical constraint maintains layer
coherence (low deviation from master) but does not create directional
preference. Emanation is thus a \emph{constrained} variant of spectral-only
training, not a topology-programming mechanism.
